% Options for packages loaded elsewhere
\PassOptionsToPackage{unicode}{hyperref}
\PassOptionsToPackage{hyphens}{url}
\documentclass[
  ukrainian,
]{article}
\usepackage{xcolor}
\usepackage[margin=2.5cm]{geometry}
\usepackage{amsmath,amssymb}
\setcounter{secnumdepth}{-\maxdimen} % remove section numbering
\usepackage{iftex}
\ifPDFTeX
  \usepackage[T1]{fontenc}
  \usepackage[utf8]{inputenc}
  \usepackage{textcomp} % provide euro and other symbols
\else % if luatex or xetex
  \usepackage{unicode-math} % this also loads fontspec
  \defaultfontfeatures{Scale=MatchLowercase}
  \defaultfontfeatures[\rmfamily]{Ligatures=TeX,Scale=1}
\fi
\usepackage{lmodern}
\ifPDFTeX\else
\setmainfont{DejaVu Serif}
\setsansfont{DejaVu Sans}
\setmonofont{DejaVu Sans Mono}
\setmathfont{STIX Two Math}
\fi
\ifPDFTeX\else
  % xetex/luatex font selection
\fi
% Use upquote if available, for straight quotes in verbatim environments
\IfFileExists{upquote.sty}{\usepackage{upquote}}{}
\IfFileExists{microtype.sty}{% use microtype if available
  \usepackage[]{microtype}
  \UseMicrotypeSet[protrusion]{basicmath} % disable protrusion for tt fonts
}{}
\makeatletter
\@ifundefined{KOMAClassName}{% if non-KOMA class
  \IfFileExists{parskip.sty}{%
    \usepackage{parskip}
  }{% else
    \setlength{\parindent}{0pt}
    \setlength{\parskip}{6pt plus 2pt minus 1pt}}
}{% if KOMA class
  \KOMAoptions{parskip=half}}
\makeatother
\usepackage{graphicx}
\graphicspath{{/Users/kalinicenkovitalijmikolajovic/.codex/worktrees/5c15/Script_for_translate_book/Output_Final/}}
\setkeys{Gin}{width=0.82\linewidth,height=0.82\textheight,keepaspectratio}
\makeatletter
\newsavebox\pandoc@box
\newcommand*\pandocbounded[1]{% scales image to fit in text height/width
  \sbox\pandoc@box{#1}%
  \Gscale@div\@tempa{\textheight}{\dimexpr\ht\pandoc@box+\dp\pandoc@box\relax}%
  \Gscale@div\@tempb{\linewidth}{\wd\pandoc@box}%
  \ifdim\@tempb\p@<\@tempa\p@\let\@tempa\@tempb\fi% select the smaller of both
  \ifdim\@tempa\p@<\p@\scalebox{\@tempa}{\usebox\pandoc@box}%
  \else\usebox{\pandoc@box}%
  \fi%
}
% Set default figure placement to htbp
\def\fps@figure{htbp}
\makeatother
\ifLuaTeX
\usepackage[bidi=basic,shorthands=off,provide=*]{babel}
\else
\usepackage[bidi=default,shorthands=off,provide=*]{babel}
\fi
\ifLuaTeX
  \usepackage{selnolig} % disable illegal ligatures
\fi
\setlength{\emergencystretch}{3em} % prevent overfull lines
\providecommand{\tightlist}{%
  \setlength{\itemsep}{0pt}\setlength{\parskip}{0pt}}
\usepackage{bookmark}
\IfFileExists{xurl.sty}{\usepackage{xurl}}{} % add URL line breaks if available
\urlstyle{same}
\hypersetup{
  pdflang={uk},
  hidelinks,
  pdfcreator={LaTeX via pandoc}}

\author{}
\date{}

\begin{document}

\pandocbounded{\includegraphics[keepaspectratio]{images/_page_359_Figure_10.jpeg}}

ATHBLK1203X

\[
w = 2t \left[ \frac{z-s}{z-t} \right],
\]

MATHBLK1204X

\[
u = \gamma [(x-s)(x-t)-\gamma^2],
\]

MATHBLK1205X

\[
v = \gamma [y(x-t)],
\]

MATHBLK1206X

\[
\gamma = 2t/[(x-t)^2 + y^2].
\]

MATHBLK1207X

\begin{verbatim}
\begin{aligned} \varphi_{uu} + \varphi_{vv} &= 0, & \text{в кольче,} \\ (\Gamma Y) & & & & & & & \\ (\Gamma Y) & & & & & & \\ (\Gamma Y) & & & & & \\ & & & & & \\ (\Gamma Y) & & & & & \\ & & & & \\ & & & & \\ & & & & \\ & & & & \\ & & & \\ & & & \\ & & & \\ & & & \\ & & & \\ & & & \\ & & \\ & & \\ & & \\ & & \\ & & \\ & & \\ & & \\ & & \\ & & \\ & & \\ & & \\ & & \\ & & \\ & & \\ & & \\ & & \\ & & \\ & & \\ & & \\ & & \\ & & \\ & & \\ & & \\ & & \\ & & \\ & & \\ & & \\ & & \\ & & \\ & & \\ & & \\ & & \\ & & \\ & & \\ & & \\ & & \\ & & \\ & & \\ & & \\ & & \\ & & \\ & & \\ & & \\ & & \\ & & \\ & & \\ & & \\ & & \\ & & \\ & & \\ & & \\ & & \\ & & \\ & & \\ & & \\ & & \\ & & \\ & & \\ & & \\ & & \\ & & \\ & & \\ & & \\ & & \\ & & \\ & & \\ & & \\ & & \\ & & \\ & & \\ & & \\ & & \\ & & \\ & & \\ & & \\ & & \\ & & \\ & & \\ & & \\ & & \\ & & \\ & & \\ & & \\ & & \\ & & \\ & & \\ & & \\ & & \\ & & \\ & & \\ & & \\ & & \\ & & \\ & & \\ & & \\ & & \\ & & \\ & & \\ & & \\ & & \\ & & \\ & & \\ & & \\ & & \\ & & \\ & & \\ & & \\ & & \\ & & \\ & & \\ & & \\ & & \\ & & \\ & & \\ & & \\ & & \\ & & \\ & & \\ & & \\ & & \\ & & \\ & & \\ & & \\ & & \\ & & \\ & & \\ & & \\ & & \\ & & \\ & & \\ & & \\ & & \\ & & \\ & & \\ & & \\ & & \\ & & \\ & & \\ & & \\ & & \\ & & \\ & & \\ & & \\ & & \\ & & \\ & & \\ & & \\ & & \\ & & \\ & & \\ & & \\ & & \\ & & \\ & & \\ & & \\ & & \\ & & \\ & & \\ & & \\ & & \\ & & \\ & & \\ & & \\ & & \\ & & \\ & & \\ & & \\ & & \\ & & \\ & & \\ & & \\ & & \\ & & \\ & & \\ & & \\ & & \\ & & \\ & & \\ & & \\ & & \\ & & \\ & & \\ & & \\ & & \\ & & \\ & & \\ & & \\ & & \\ & & \\ & & \\ & & \\ & & \\ & & \\ & & \\ & & \\ & & \\ & & \\ & & \\ & & \\ & & \\ & & \\ & & \\ & & \\ & & \\ & & \\ & & \\ & & \\ & & \\ & & \\ & & \\ & & \\ & & \\ & & \\ & & \\ & & \\ & & \\ & & \\ & & \\ & & \\ & & \\ & & \\ & & \\ & & \\ & & \\ & & \\ & & \\ & & \\ & & \\ & & \\ & & \\ & & \\ & & \\ & & \\ & & \\ & & \\ & & \\ & & \\ & & \\ & & \\ & & \\ & & \\ & & \\ & & \\ & & \\ & & \\ & & \\ & & \\ & & \\ & & \\ & & \\ & & \\ & & \\ & & \\ & & \\ & & \\ & & \\ & & \\ & & \\ & & \\ & & \\ & & \\ & & \\ & & \\ & & \\ & & \\ & & \\ & & \\ & & \\ & & \\ & & \\ & & \\ & & \\ & & \\ & & \\ & & \\ & & \\ & & \\ & & \\ & & \\ & & \\ & & \\ & & \\ & & \\ & & \\ & & \\ & & \\ & & \\ & & \\ & & \\ & & \\ & & \\ & & \\ & & \\ & & \\ & & \\ & & \\ & & \\ & & \\ & & \\ & & \\ & & \\ & & \\ & &
\end{verbatim}

MATHBLK1208X

\[
\begin{cases} 4/n\pi, & n = 1, 3, \dots, \\ 0, & n = 2, 4, \dots, \end{cases}
\]

MATHBLK1209X

\[
\varphi(r) = a \ln r + b, \quad r = u^2 + v^2.
\]

MATHBLK1210X

\[
\varphi(u, v) = 0.57 \ln(u^2 + v^2) + 1
\]

MATHBLK1211X

\[
\varphi(x, y) = 0.57 \ln(u^2 + v^2) + 1,
\]

MATHBLK1212X

\[
w = \ln \left[ \frac{z-1}{z+1} \right].
\]

MATHBLK1213X

\[
z = re^{i\theta}
\]

MATHBLK1214X

\[
\begin{cases} 1, & n = 1, \\ 0, & n = 2, 3, \dots, \end{cases}
\]

MATHBLK1215X

\begin{verbatim}
\begin{aligned} (\text{Y} \forall \Pi) & \qquad \phi_{\text{vx}} + \phi_{yy} = 0, \quad 0 < x < \infty, \quad 0 < y < \infty, \\ \left\{ \begin{array}{l} \phi(x, \ 0) = \begin{cases} 1, & 0 < x < 1, \\ 0, & 1 \leqslant x < \infty, \end{cases} \\ \phi(0, \ y) = \begin{cases} 1, & 0 < y < 1, \\ 0, & 1 \leqslant y < \infty. \end{cases} \end{aligned}
\end{verbatim}

MATHBLK1216X

MATHBLK1217X

MATHBLK1218X

MATHBLK1219X

\begin{verbatim}
\begin{cases} S_n(0) = \begin{cases} 4/n\pi, & n = 1, 3, \dots, \\ 0, & n = 2, 4, \dots, \end{cases}
\end{verbatim}

MATHBLK1220X

MATHBLK1221X

MATHBLK1222X

MATHBLK1223X

MATHBLK1224X

\[
\begin{cases} \frac{dS_n(0)}{dt} = 0. \end{cases}
\]

MATHBLK1225X

MATHBLK1226X

MATHBLK1227X

MATHBLK1228X,

где

MATHBLK1229X

MATHBLK1230X

MATHBLK1231X

MATHBLK1232X

MATHBLK1233X

MATHBLK1234X

MATHBLK1235X

\begin{enumerate}
\def\labelenumi{\arabic{enumi}.}
\setcounter{enumi}{8}
\tightlist
\item
\end{enumerate}

MATHBLK1236X

MATHBLK1237X

\[
\begin{cases} u(0, t) = f(t) \\ u(L, t) = g(t) \end{cases}
\]

MATHBLK1238X

MATHBLK1239X

\(f(x)\) * \(g(x)\)

MATHBLK1240X

MATHBLK1241X

MATHBLK1242X

\begin{enumerate}
\def\labelenumi{\arabic{enumi}.}
\setcounter{enumi}{1}
\tightlist
\item
  Решите в общем виде задачу
\end{enumerate}

MATHBLK1243X

MATHBLK1244X

MATHBLK1245X

MATHBLK1246X

MATHBLK1247X

\section{ПРИНЦИП СУПЕРПОЗИЦІЇ Є ОСНОВОЮ ТЕОРІЇ ЛІНІЙНИХ
СИСТЕМ}\label{ux43fux440ux438ux43dux446ux438ux43f-ux441ux443ux43fux435ux440ux43fux43eux437ux438ux446ux456ux457-ux454-ux43eux441ux43dux43eux432ux43eux44e-ux442ux435ux43eux440ux456ux457-ux43bux456ux43dux456ux439ux43dux438ux445-ux441ux438ux441ux442ux435ux43c}

МЕТА ЛЕКЦІЇ: Ознайомити читача з принципом суперпозиції та показати, як
його можна використати для розкладання початкової задачі на підзадачі,
потім розв'язання всіх підзадач одночасно і, підсумувавши отримані
результати, знаходження рішення початкової задачі. Ми також хочемо
показати, що два основні методи розв'язання лінійних рівнянь --- метод
розділення змінних і метод інтегрального перетворення --- використовують
принцип суперпозиції.

Припустимо, що інженеру потрібно знайти відгук лінійної системи на
вхідний ефект f.~Загальний підхід до розв'язання задачі такий:

\begin{itemize}
\tightlist
\item
  \begin{enumerate}
  \def\labelenumi{\arabic{enumi})}
  \tightlist
  \item
    розкласти f на елементарні ефекти \(f = \sum f_k\) ;
  \end{enumerate}
\item
  \begin{enumerate}
  \def\labelenumi{\arabic{enumi})}
  \setcounter{enumi}{1}
  \tightlist
  \item
    знайти реакцію системи \(u_k\) на вплив \(f_k\) ;
  \end{enumerate}
\end{itemize}

\begin{enumerate}
\def\labelenumi{\arabic{enumi})}
\setcounter{enumi}{2}
\tightlist
\item
  сложить отклики \(u = \sum_{i=1}^{n} u_{i}\) .
\end{enumerate}

Якщо система лінійна, то отримана сума є реакцією системи на вхідний
ефект f.~Це принцип суперпоезії (див. стор. 26.1).

MATHBLK1248X

РИС. 26.1. Основна ідея принципу суперпоенції; L --- диференціальний
оператор; f --- вхід; У---виходь.

Ця базова ідея може використовуватися для розв'язання складних змішаних
задач, розбиваючи їх на прості підзавдання, які розв'язуються за
допомогою

кожне з них, а потім підсумовування отриманих результатів (звісно,
диференціальне рівняння і граничні умови мають бути лінійними).

\section{Розкладання змішаної задачі на дві
простіші}\label{ux440ux43eux437ux43aux43bux430ux434ux430ux43dux43dux44f-ux437ux43cux456ux448ux430ux43dux43eux457-ux437ux430ux434ux430ux447ux456-ux43dux430-ux434ux432ux456-ux43fux440ux43eux441ux442ux456ux448ux456}

Предположим, нам необходимо решить линейную задачу (обозначим ее P)

(УЧП)

MATHBLK1249X

(Р) (ГУ)

MATHBLK1250X

MATHBLK1251X

\[
\begin{cases} u(0, t) = 0, \\ u(1, t) = 0, \\ u(x, 0) = \sin(2\pi x), \quad 0 \le x \le 1. \end{cases}
\]

MATHBLK1252X

MATHBLK1253X

MATHBLK1254X

MATHBLK1255X

MATHBLK1256X

MATHBLK1257X

MATHBLK1258X

\[
\begin{cases} u(0, t) = 0, \\ u(1, t) = 0, \\ u(1, t) = 0, \end{cases}
\]

MATHBLK1259X

MATHBLK1260X

MATHBLK1261X

MATHBLK1262X

MATHBLK1263X

MATHBLK1264X

MATHBLK1265X

MATHBLK1266X

MATHBLK1267X

MATHBLK1268X

MATHBLK1269X

\[
\begin{cases} u_t = u_{xx} + f(x, t), \\ u(0, t) = 0, \\ u(1, t) = 0, \\ u(x, 0) = \varphi(x) \end{cases}
\]

MATHBLK1270X

MATHBLK1271X

\[
\cos(ax)
\]

MATHBLK1272X

MATHBLK1273X

\[
\begin{cases} u_t = u_{xx} + f(x, t), \\ u(0, t) = 0, \\ u(1, t) = 0, \\ u(x, 0) = 0 \end{cases}
\]

MATHBLK1274X

MATHBLK1275X

MATHBLK1276X

MATHBLK1277X

MATHBLK1278X

\[
\begin{cases} u_t = u_{xx}, \\ u(0, t) = 0, \\ u(1, t) = 0, \\ u(x, 0) = \varphi(x). \end{cases}
\]

MATHBLK1279X

MATHBLK1280X

f'\,'(0)

MATHBLK1281X

MATHBLK1282X

MATHBLK1283X

MATHBLK1284X

MATHBLK1285X

MATHBLK1286X

MATHBLK1287X

MATHBLK1288X

MATHBLK1289X

MATHBLK1290X

MATHBLK1291X

MATHBLK1292X

MATHBLK1293X

MATHBLK1294X

\[
\begin{split} A_n(t) &= 2\int_0^1 u_t \sin(n\pi x) \, dx = \frac{d}{dt} \bigg[ 2\int_0^1 u(x, t) \sin(n\pi x) \, dx \bigg] = \frac{dU_n(t)}{dt} \,, \\ B_n(t) &= 2\int_0^1 u_{xx} \sin(n\pi x) \, dx = -(n\pi)^2 \, U_n(t) \,+ \\ &\quad + 2n\pi \, \big[ u(0, t) + (-1)^{n+1} \, u(1, t) \big], \end{split}
\]

MATHBLK1295X

MATHBLK1296X

MATHBLK1297X

MATHBLK1298X

MATHBLK1299X

\[
\begin{cases} u(0, t) = 0, \\ u(1, t) = 0 \end{cases}
\]

MATHBLK1300X

MATHBLK1301X

MATHBLK1302X

MATHBLK1303X

MATHBLK1304X

MATHBLK1305X

MATHBLK1306X

MATHBLK1307X

MATHBLK1308X

MATHBLK1309X

MATHBLK1310X

MATHBLK1311X

\[
0 \le x \le \pi
\]

MATHBLK1312X

MATHBLK1313X

MATHBLK1314X

f'\,'(x)

MATHBLK1315X

MATHBLK1316X

MATHBLK1317X

MATHBLK1318X

MATHBLK1319X

MATHBLK1320X

MATHBLK1321X

MATHBLK1322X

MATHBLK1323X

MATHBLK1324X

\begin{verbatim}
\begin{array}{ll} (\mathrm{Y} + \Pi) & u_{tt} = u_{xx} + \sin{(3\pi x)}, & 0 < x < 1, & 0 < t < \infty, \\ (\mathrm{\Gamma} \mathrm{Y}) & \left\{ \begin{array}{ll} u\left(0, \ t\right) = 0, \\ u\left(1, \ t\right) = 0, \\ \end{array}
\end{verbatim}

MATHBLK1325X

MATHBLK1326X

MATHBLK1327X

MATHBLK1328X

Для каждой из подзадач выберите метод решения по своему BKYCY,

\begin{enumerate}
\def\labelenumi{\arabic{enumi}.}
\setcounter{enumi}{2}
\tightlist
\item
  Нехай \(u_1\) і \(u_2\) є розв'язками наступних рівнянь:
\end{enumerate}

MATHBLK1329X;\\
b) \(u_t = u_{xx} + e^t\) ;

\begin{enumerate}
\def\labelenumi{\Alph{enumi})}
\setcounter{enumi}{1}
\tightlist
\item
\end{enumerate}

MATHBLK1330X;\\
F) \(u_t = u_{xx} + u^2\) .

MATHBLK1331X

MATHBLK1332X

\[
\begin{aligned} u_t &= u_{xx} + f(x, t), & 0 < x < 1, & 0 < t < \infty, \\ & \begin{cases} u(0, t) = g_1(t), \\ u(1, t) = g_2(t), \end{cases} & 0 < t < \infty, \\ & u(x, 0) = \varphi(x), & 0 \le x \le 1. \end{aligned}
\]

MATHBLK1333X

MATHBLK1334X

x

MATHBLK1335X

MATHBLK1336X

\_ \textbar, -\infty \textless{} x \textless{} \infty  , 0 \textless{} t
\textless{} \infty  ,

MATHBLK1337X

\[
(HV)   u(x, 0) = \varphi(x), \quad -\infty < x < \infty.
\]

MATHBLK1338X

\[
\Delta u = u_{xx} + u_{yy},   -\infty < x < \infty  ,   0 < t < \infty
\]

MATHBLK1339X

\[
\Delta u = u_{rr} + \frac{1}{r}u_r + \frac{1}{r^2}u_{\theta\theta}
\]

MATHBLK1340X

\[
\Delta u = u_{xx} + u_{yy} + u_{zz}
\]

MATHBLK1341X

\[
\Delta u = u_{rr} + \frac{1}{r}u_r + \frac{1}{r^2}u_{\theta\theta} + u_{zz}.
\]

MATHBLK1342X M A

\pandocbounded{\includegraphics[keepaspectratio]{images/_page_196_Figure_10.jpeg}}

THBLK1343X

\[
\Delta u = u_{rr} + \frac{2}{r}u_r + \frac{1}{r^2}u_{\theta\theta} + \frac{\cot\theta}{r^2}
\]

MATHBLK1344X M A

\pandocbounded{\includegraphics[keepaspectratio]{images/_page_197_Figure_5.jpeg}}

THBLK1345X

\[
\Delta u = u_{rr} + \frac{2}{r}u_r + \frac{1}{r^2}u_{\theta\theta} + \frac{\cot\theta}{r^2}u_0 + \frac{1}{r^2\sin\theta}u_{\phi\phi}
\]

MATHBLK1346X M

\end{document}
